\section{The Polygonal Model}
\label{sec:polygonal-model}

\begin{layman}
There is a classical trick — known to topologists since the late
nineteenth century — for turning a curved closed surface into a
single flat sheet of paper. Take a doughnut. Cut along the loop
that goes through the hole; you've turned it into a tube. Cut the
tube once around its length; it unrolls into a flat rectangle.
The four edges of the rectangle, suitably glued back together,
recover the doughnut. Higher-genus surfaces work the same way: a
two-handled pretzel cuts open into an octagon, a three-handled one
into a twelve-gon, and a genus-\(g\) surface into a \(4g\)-gon.
The boundary identifications follow a single canonical pattern:
\(a_1 b_1 a_1^{-1} b_1^{-1} \cdots a_g b_g a_g^{-1} b_g^{-1}\),
four arcs per handle.

This is the \emph{polygonal model}, and it is the bridge between
two languages. On the original surface, integration is hard: there
are loops that nothing can fill in, so closed forms can integrate
to nonzero numbers around those loops, and Stokes' theorem doesn't
apply directly. On the unfolded polygon, the interior is a flat
disk — simply connected, no loops at all — and Stokes becomes
elementary. Every period integral on the surface translates into a
boundary integral on the polygon, and the canonical gluing pattern
lets us pair up boundary edges to recover the original loops.
The \emph{Riemann bilinear relations} of the next section are
exactly the residue of this translation.

The polygonal model also bridges \emph{analytic} and
\emph{topological} genus. The analytic genus we built in section
2 is a complex vector-space dimension; the topological genus is
``half the rank of \(H_1(X,\Z)\)'' — pure topology, blind to the
complex structure. The polygon's \(2g\) symplectic loops give an
explicit basis of homology, computing the topological genus
directly; the polygonal homeomorphism then transports that
computation back to the original surface, finally identifying the
two notions of genus. Without this bridge, the genus-zero
classification of section 4 would be circular.

Two technical projects sit inside this chapter. \textbf{Stage A},
the \emph{surface classification theorem}, is the homeomorphism
\(X \cong \mathrm{Polygon4g}(g)\) itself. Buzzard's challenge cannot
work without it, but Mathlib does not yet contain it. The current
formalization develops the combinatorial side — edge words, the
standard relator, Tietze moves — as the substrate on which a future
proof will rest. \textbf{Stage B}, the \emph{topological-genus
bridge}, is the chain that takes a complex curve, derives its real
two-dimensional manifold structure, computes its first homology
group via singular chains, and divides by two to recover the
analytic genus. A non-trivial fraction of Stage B is already
sorry-free; the remainder is on a known path.
\end{layman}

A compact connected oriented Riemann surface \(X\) of genus \(g\) is
homeomorphic to a single flat \(4g\)-gon \(P\) with edges glued back
together in the canonical pattern
\(a_1b_1a_1^{-1}b_1^{-1}\cdots a_gb_ga_g^{-1}b_g^{-1}\). This chapter
develops that polygonal model as a self-contained subsystem: the
quotient construction \(\mathrm{Polygon4g}\), the edge-word algebra
that certifies the standard boundary identification, the smooth real
2-manifold structure underlying any complex curve, the topological
genus, and the analysis-on-the-polygon lemmas. The next chapter then
builds the Riemann bilinear relations on top of this scaffold, by
running Stokes' theorem on the open polygon and folding edges
together pairwise via the identification pattern.

\subsection{The fundamental \(4g\)-gon}

\begin{definition}[Boundary parameterization of the closed unit disk]
\label{def:polygon4g-boundary-param}
For each genus \(g\geq 1\) and each index \(i\in\{0,\ldots,4g-1\}\),
the \(i\)-th boundary arc of the disk is the map
\(t\mapsto e^{i\,2\pi(i+t)/4g}\) for \(t\in[0,1]\), lifted into the
closed unit disk. The \(4g\) arcs cover the boundary circle once
each.
\lean{JacobianChallenge.Periods.boundaryParamC}
\leanok
\end{definition}

\begin{layman}
Take the closed unit disk in the plane. March around its boundary
circle and chop the lap into \(4g\) equal arcs, numbered
\(0,1,\ldots,4g-1\). Each arc gets a parameter \(t\in[0,1]\)
running from start to end. This is just preparation: every face of
the future polygon will turn out to be one of these arcs, and the
gluing rule will be expressed as a map between paired arcs. Why
bother? It puts a concrete, computable structure onto ``the
boundary of the polygon,'' so that the later equivalence relation
can be defined by a simple formula in \(t\) rather than an abstract
combinatorial chart.
\end{layman}

\begin{definition}[Standard \(4g\)-gon side-pairing relation]
\label{def:polygon4g-side-rel}
\uses{def:polygon4g-boundary-param}
The \emph{standard side-pairing relation} on the closed unit disk
identifies, for every handle index \(i\in\{0,\ldots,g-1\}\) and every
parameter \(t\in[0,1]\):
\begin{itemize}
  \item the point on arc \(4i\) at parameter \(t\) with the point on
    arc \(4i+2\) at parameter \(1-t\) (the \(a_i\)-cycle);
  \item the point on arc \(4i+1\) at parameter \(t\) with the point on
    arc \(4i+3\) at parameter \(1-t\) (the \(b_i\)-cycle).
\end{itemize}
The full relation is the equivalence closure of these generating
pairs.
\lean{JacobianChallenge.Periods.Polygon4g.SideRel}
\leanok
\end{definition}

\begin{layman}
Pair up the \(4g\) arcs in groups of four — one group per handle.
Within each group of four consecutive arcs, glue the first to the
third with reversed orientation, and the second to the fourth with
reversed orientation. ``With reversed orientation'' means that the
parameter \(t\) running from \(0\) to \(1\) on one arc gets matched
against the parameter running from \(1\) to \(0\) on its partner.
For a torus (\(g=1\)) the pattern matches the top edge to the bottom
and the left to the right of a square — the classical wrap-around.
For higher genus, every handle adds another four-arc block. Why
bother? This single algebraic rule — applied uniformly for every
handle — is what turns the disk into a higher-genus surface after
quotienting.
\end{layman}

\begin{definition}[Standard fundamental polygon \(\mathrm{Polygon4g}\)]
\label{def:polygon4g}
\uses{def:polygon4g-side-rel}
The standard fundamental \(4g\)-gon is the topological quotient of
the closed unit disk by the standard side-pairing relation:
\[
  \mathrm{Polygon4g}(g) := \overline{D}^{1}\,/\!\sim_{\mathrm{side}}.
\]
Its quotient map is denoted \(\mathrm{mk}_g\) and is continuous by
construction.
\lean{JacobianChallenge.Periods.Polygon4g}
\leanok
\end{definition}

\begin{layman}
Take the closed unit disk and collapse together every pair of
boundary points that the side-pairing relation declared equivalent.
The resulting space is the standard genus-\(g\) fundamental polygon.
Topologically, the boundary disappears (it has been folded onto
itself) and the quotient is a compact surface — for \(g=1\), an
ordinary torus; for \(g=2\), a two-handled pretzel; and so on. The
quotient map sends each disk point to its equivalence class, and is
continuous because the equivalence relation is generated by a closed
set of identifications. Why bother? Every topological subtlety in
the surface — its handles, its homology, its non-simply-connectedness
— is now encoded in the boundary identifications, while the
\emph{interior} of the disk remains a flat, simply connected piece
of paper. That separation is what makes Stokes-on-the-polygon
tractable.
\end{layman}

\subsection{Edge-word combinatorics}

The polygonal model is most naturally manipulated through its
\emph{boundary word}: a sequence of letters from a four-letter
alphabet that records, in order, the identification class of each
boundary arc. The standard side-pairing corresponds to a specific
word; alternative words give homeomorphic surfaces if and only if
they are related by an explicit set of combinatorial moves (Tietze
transformations).

\begin{definition}[Edge-word alphabet]
\label{def:edge-word-alphabet}
For each genus \(g\), the edge-word alphabet \(\mathrm{Letter}(g)\)
has four families of letters indexed by handle \(i\in\{0,\ldots,g-1\}\):
\(a_i\), \(b_i\), \(a_i^{-1}\), \(b_i^{-1}\).
An \emph{edge word} is a finite list of letters,
\(\mathrm{EdgeWord}(g) := \mathrm{List}(\mathrm{Letter}(g))\).
\lean{JacobianChallenge.Periods.Letter, JacobianChallenge.Periods.EdgeWord}
\leanok
\end{definition}

\begin{layman}
Pick a four-letter alphabet for each handle: ``go around the
\(a\)-loop,'' ``go around the \(b\)-loop,'' and their inverses ``go
backwards around \(a\),'' ``go backwards around \(b\).'' An edge
word is a finite sequence of such letters — a walking direction
through the boundary loops. For a single-handle surface (genus 1)
the alphabet has four letters and a typical word is \(aba^{-1}b^{-1}\).
Why bother? Storing the gluing pattern as a list of letters lets us
manipulate it combinatorially: we can ask when two patterns produce
homeomorphic surfaces, decompose handles, cancel inverses, and run
the surface-classification algorithm — all as list operations.
\end{layman}

\begin{definition}[Standard relator word]
\label{def:standard-word}
\uses{def:edge-word-alphabet}
The \emph{standard relator word} of genus \(g\) is the concatenation
\[
  \mathrm{standardWord}(g)
    := a_0\,b_0\,a_0^{-1}\,b_0^{-1}\,
       a_1\,b_1\,a_1^{-1}\,b_1^{-1}\,\cdots\,
       a_{g-1}\,b_{g-1}\,a_{g-1}^{-1}\,b_{g-1}^{-1}.
\]
A word is in \emph{standard form} if it equals
\(\mathrm{standardWord}(g)\).
\lean{JacobianChallenge.Periods.EdgeWord.standardWord}
\leanok
\end{definition}

\begin{layman}
Concatenating the four-letter handle blocks for every
\(i=0,1,\ldots,g-1\) gives one canonical word of length \(4g\). For
\(g=2\) the standard relator is \(a_0 b_0 a_0^{-1} b_0^{-1}
a_1 b_1 a_1^{-1} b_1^{-1}\). This is the precise word every compact
orientable genus-\(g\) surface ``simplifies to'' under the topological
classification theorem. Why bother? Naming this canonical word fixes
a normal form. Every other gluing pattern can be checked against
this template; when an arbitrary edge word is shown equivalent
(combinatorially) to the standard one, that's a certificate that the
underlying surface is genuinely the genus-\(g\) closed orientable
surface and not some exotic close cousin.
\end{layman}

\begin{lemma}[Standard word's quotient is \(\mathrm{Polygon4g}\)]
\label{lem:polygon4g-eq-standard-word-quotient}
\uses{def:polygon4g,def:standard-word}
The quotient of the closed unit disk by the side-pairing induced by
\(\mathrm{standardWord}(g)\) is identical (as a type) to
\(\mathrm{Polygon4g}(g)\).
\lean{JacobianChallenge.Periods.EdgeWord.polygon4g_eq_standard_word_quotient}
\leanok
\end{lemma}
\begin{proof}\leanok\end{proof}

\begin{layman}
Two routes to ``the polygon'' coincide. One: directly identify the
\(4g\) boundary arcs by the standard pattern and quotient. Two:
write down the standard relator word, derive its arc-pairing rule
from the word's letters, and quotient. The lemma says these two
constructions land in the same type. Why bother? It lets us freely
swap between the geometric construction (used for topology and
analysis on the polygon) and the combinatorial one (used for surface
classification and Tietze moves), without translating back and forth
through an awkward homeomorphism.
\end{layman}

\begin{definition}[Tietze equivalence on edge words]
\label{def:tietze-equivalence}
\uses{def:edge-word-alphabet}
Two edge words are \emph{Tietze-equivalent} if one can be obtained
from the other by a finite sequence of moves drawn from:
\begin{itemize}
  \item \emph{Inverse cancellation}: insert or remove a consecutive
    pair \(\ell\,\ell^{-1}\) anywhere in the word;
  \item \emph{Handle swap}: permute two complete four-letter handle
    blocks within the word.
\end{itemize}
The relation is generated as the equivalence closure of these single
moves.
\lean{JacobianChallenge.Periods.EdgeWord.TietzeEq}
\leanok
\end{definition}

\begin{layman}
Two recipes for gluing the boundary of a polygon are
\emph{Tietze-equivalent} if you can transform one into the other by
a sequence of small, combinatorially trivial steps: cancel a
generator immediately followed by its inverse (because gluing
\(\ell\) and then ungluing \(\ell\) is the same as not gluing); or
swap the order of two complete four-letter handle blocks (because
relabeling which handle is ``first'' doesn't change the surface).
Why bother? Tietze equivalence is the combinatorial side of
``producing the same topological surface.'' If we can show every
\emph{some} edge word for a given surface is Tietze-equivalent to the
standard word, we get a normal-form theorem — and that's exactly the
heart of the surface-classification step (Stage A) below.
\end{layman}

\subsection{Classical inputs for surface classification}

\begin{theorem}[Classical input: Radó's triangulation theorem]
\label{input:rado-triangulation}
\uses{lem:rado-pl-atlas-exists,lem:rado-assembled-simplicial-complex,lem:rado-realisation-homeomorphism,lem:rado-overview-stepwise,lem:rado-pl-assembled-is-2mfd,lem:rado-pl-assembled-realisation-homeo,lem:rado-overview-glued-continuous,lem:rado-overview-glued-bijective}
Every compact connected oriented topological 2-manifold admits a
triangulation: a homeomorphism with the underlying space of a finite
two-dimensional simplicial complex.
\depends{Classical 1925 result (T. Radó, \emph{Über den Begriff der
Riemannschen Fläche}). Mathlib v4.28.0 has neither the
simplicial-complex side nor the existence theorem; the formalisation
is a multi-month project in its own right, conventionally treated
through PL or smooth-triangulation methods. The project carries a
bottom-up sketch of the classical proof in
\code{Jacobian/StageA/RadoTheorem.lean} (chart cover, PL
approximation, combinatorial assembly, ${\sim}150$ LOC of typed
sorries) backed by the supporting modules
\code{StageA/SimplicialComplex.lean} and
\code{StageA/SpanningTree.lean}.}
\lean{JacobianChallenge.Blueprint.Sec03.rado_triangulation}
\notready
\end{theorem}

\begin{layman}
Radó's theorem says every compact two-dimensional surface — Riemann
surface or not — can be cut into a finite collection of triangles
glued along their edges, in a way that respects the topology
exactly. For a sphere, three or four triangles suffice. For a
torus, two triangles glued into a square (with the canonical wrap-
around) work. For higher-genus surfaces, more triangles are needed,
but always finitely many. Why bother? The polygonal-model theorem
runs on triangulations: cut the triangulation into a tree, lay it
flat as a polygon, run Tietze moves to reduce the boundary word.
Without Radó's theorem the entire Stage-A machinery has no input.
Mathlib does not yet have it; the project carries a thin
\code{rado\_triangulation} stub naming the classical result, to be
filled by a future formalisation.
\end{layman}

\begin{theorem}[Classical input: Tietze normal form for orientable surface words]
\label{input:tietze-normal-form}
\uses{def:tietze-equivalence,def:orientable,lem:tietze-to-handle-grouped,lem:tietze-handle-grouped-to-standard,lem:tietze-overview-stepwise}
Every edge word arising from a polygonal presentation of a compact
connected orientable 2-manifold of genus \(g\) is Tietze-equivalent to
the standard relator
\(a_1 b_1 a_1^{-1} b_1^{-1} \cdots a_g b_g a_g^{-1} b_g^{-1}\).
\depends{Classical Brahana--Seifert--Threlfall reduction. Mathlib
v4.28.0 has neither edge words nor Tietze transformations on them;
the project ships a top-down sketch under
\code{Jacobian/StageA/EdgeWordTietze.lean}
(\code{orientable_edgeWord_tietzeEq_standardWord}) and a parallel
project-side decomposition in \code{Jacobian/Periods/TietzeReduction.lean}
(handle separation, inverse cancellation, cyclic reduction).}
\lean{JacobianChallenge.StageA.orientable_edgeWord_tietzeEq_standardWord}
\leanok
\end{theorem}

\begin{layman}
The Tietze normal-form theorem is the combinatorial heart of
surface classification. Take any cell-complex presentation of an
orientable closed surface: triangulate it, cut along a spanning
tree, lay it flat as a polygon, and read off the boundary word — a
sequence in the four letters \(a_i, b_i, a_i^{-1}, b_i^{-1}\). The
theorem says: by a finite sequence of \emph{Tietze moves} (cancel
adjacent inverse pairs; permute complete handle blocks), every such
word transforms into the standard relator
\(a_1 b_1 a_1^{-1} b_1^{-1} \cdots a_g b_g a_g^{-1} b_g^{-1}\).
For a torus (\(g = 1\)) the standard form is \(aba^{-1}b^{-1}\); for
a genus-2 surface, \(a_1 b_1 a_1^{-1} b_1^{-1} a_2 b_2 a_2^{-1}
b_2^{-1}\); and so on. Why bother? Combined with Radó's theorem
(triangulation always exists), Tietze normal form gives the
polygonal-model theorem: every orientable surface looks
topologically like the standard \(4g\)-gon. The Lean side
formalises the move-by-move reduction algorithm; the headline
theorem currently sits at the top of a sorry tree.
\end{layman}

\subsection{The polygonal model theorem (Stage A)}

\begin{theorem}[Polygonal model of a compact Riemann surface]
\label{thm:polygonal-model}
\uses{def:polygon4g,lem:polygon4g-eq-standard-word-quotient,def:tietze-equivalence,def:orientable,input:rado-triangulation,input:tietze-normal-form,lem:polygonal-model-some-polygon-homeomorph,lem:polygonal-model-genus-matches,lem:polygonal-model-overview-stepwise,lem:polygonal-model-phase12-word-quotient,lem:polygonal-model-phase3-tietze-normal-form,lem:polygonal-model-phase3-quotient-lift,lem:polygonal-model-apply-tietze}
A compact connected oriented Riemann surface \(X\) of analytic
genus \(g\) is homeomorphic to \(\mathrm{Polygon4g}(g)\):
\[
  X \cong_{\mathrm{homeo}} \mathrm{Polygon4g}(g).
\]
\depends{Standard topological classification of compact orientable
surfaces (triangulation + reduction to the standard relator word via
Tietze moves), or a Morse-theoretic handle decomposition. Mathlib
v4.28.0 has none of these as packaged theorems; this is the largest
classical-topology gap on the polygonal-model route. Tracked
internally as ``Stage A,'' a roughly 3,000--5,000 LOC formalization
project in its own right. The project ships a top-down sketch of
the proof tree under \code{Jacobian/StageA/} (seven files: simplicial
complexes, Radó, spanning tree, prism operator, Hurewicz, edge-word
Tietze, cellular--singular comparison; ${\sim}120$ typed declarations,
all sorry) and a parallel project-side chain in
\code{Jacobian/Periods/SurfaceClassification.lean} +
\code{TietzeReduction.lean} + \code{Polygon4gCellular.lean} +
\code{H1EvenBasisViaSurfaceClassification.lean} that wires the
classical pieces into the analytic genus equality the rest of the
blueprint consumes. Edge-word algebra and the Tietze relation are
sorry-free; the umbrella reduction theorem and the Stage~A bottom-up
infrastructure remain open.}
\lean{JacobianChallenge.Blueprint.polygonal_model}
\leanok
\end{theorem}

\begin{layman}
Every compact connected oriented genus-\(g\) Riemann surface can be
``unfolded'' onto a single flat \(4g\)-gon with the standard boundary
identification pattern. The unfolding is a homeomorphism: the surface
and the polygon are topologically the same. For a torus (\(g=1\))
the polygon is a square; for genus \(2\), an octagon; for genus
\(g\), a \(4g\)-gon. Inside the polygon there are no loops at all —
it's a flat, simply connected disk. All the topological complexity
of the surface lives in the gluing pattern on the boundary. Why
bother? It converts a global, curved object into a local, flat one.
On a flat polygon, Stokes' theorem and primitive constructions
become pedestrian; the period-lattice argument and the Riemann
bilinear relations all run on this scaffolding. The proof itself
goes through ``triangulate the surface, reduce the resulting boundary
word to the standard form via Tietze moves'' — the standard
classification of compact orientable surfaces, formalized over a
multi-thousand-line stretch of project Stage A.
\end{layman}

\begin{proofsketch}
Triangulate \(X\) (Radó's theorem: every compact surface admits a
triangulation). Cut along a maximal tree of edges to lay the
triangulation flat as a polygon. Read off the resulting boundary
word, an alternating-sign list of edge labels. Apply Tietze moves
(\ref{def:tietze-equivalence}) to reduce the word to the standard
relator. The induced map on the disk descends to a homeomorphism
between \(X\) and \(\mathrm{Polygon4g}(g)\) by
Lemma~\ref{lem:polygon4g-eq-standard-word-quotient}.
\end{proofsketch}

\subsection{Smooth and topological structure (Stage B)}

The polygonal-model theorem above is purely topological: it produces
a homeomorphism, not a diffeomorphism. To bridge analytic genus
(dimension of holomorphic 1-forms) with topological genus (rank of
\(H_1(X,\Z)\) divided by two), we record the smooth-real-2-manifold
structure underlying any complex curve and the singular-homology
definition of topological genus.

\begin{definition}[Smooth real 2-manifold structure on a complex curve]
\label{def:smooth-real-structure}
A complex one-dimensional manifold \(X\) carries an induced smooth
real 2-dimensional manifold structure, obtained by composing each
complex chart with the real-coordinate equivalence
\(\C \xrightarrow{\sim} \R\times\R \xrightarrow{\sim}
\mathrm{EuclideanSpace}\,\R\,(\mathrm{Fin}\,2)\).
\lean{JacobianChallenge.Periods.complexChartedSpaceReal2,
JacobianChallenge.Periods.ChartedSpaceComplex_to_smoothReal2}
\leanok
\end{definition}

\begin{layman}
A complex one-dimensional manifold is the same data as a smooth
real two-dimensional manifold with extra structure: the complex
multiplication on each tangent space picks out a preferred
\(90^{\circ}\)-rotation, but the underlying smooth atlas is just
charts to \(\R^2\). Concretely, every complex chart \(z=x+iy\) is
secretly a real chart \((x,y)\), and complex differentiability
implies real smoothness. The construction here packages that
identification once and for all, so downstream code can talk about
``the underlying real surface'' without re-deriving charts. Why
bother? Topological genus is defined on real 2-manifolds via
singular homology; analytic genus lives on complex curves. The
bridge between them runs through this smooth-real-2-manifold
structure.
\end{layman}

\begin{definition}[Topological genus]
\label{def:topological-genus}
\uses{def:smooth-real-structure}
The \emph{topological genus} of a compact connected topological
space \(M\) is half the \(\Z\)-rank of its first singular homology:
\[
  g_{\mathrm{top}}(M) := \tfrac{1}{2}\,\mathrm{rank}_{\Z}\,H_1(M;\Z).
\]
For compact connected oriented 2-manifolds the value equals the
classical genus.
\lean{JacobianChallenge.Periods.topologicalGenus,
JacobianChallenge.Periods.singularH1}
\leanok
\end{definition}

\begin{layman}
The \emph{topological genus} of a compact connected space is half
the rank of its integer first homology group. For a compact
orientable surface of genus \(g\), the rank of \(H_1\) is \(2g\)
(the symplectic loops contribute two generators per handle), so
half of that is exactly \(g\). For a sphere, \(H_1 = 0\) and the
genus is \(0\). For a torus, \(H_1 = \Z^2\) and the genus is \(1\).
Why bother? It expresses the topological notion of genus in a
purely homological form, ready to be compared against the analytic
notion (dimension of holomorphic 1-forms) by combining with the
polygonal-model theorem and the period-lattice rank computation.
\end{layman}

\begin{theorem}[Cellular homology computes \(H_1\) of the fundamental polygon]
\label{thm:polygon4g-cellular-singular}
\uses{def:polygon4g,def:topological-genus}
For the standard \(4g\)-gon, the free abelian group generated by the
\(2g\) edge cycles is linearly isomorphic to singular first homology:
\[
  \Z^{2g} \cong H_1(\mathrm{Polygon4g}(g);\Z).
\]
\depends{The Lean assembly is now split across
\code{Jacobian/Periods/CellularChainComplex.lean},
\code{Jacobian/Periods/Hurewicz.lean}, and
\code{Jacobian/Periods/Polygon4gCellular.lean}.  The top-level theorem
\code{polygon4g\_cellular\_singular\_iso} is sorry-free: it composes
\code{polygon4g\_cellularH1\_free\_assembly} with
\code{polygon4g\_cellularH1\_to\_singularH1}.  The remaining named
leaves are (i) existence of the standard cellular model on the quotient
polygon, (ii) the cellular boundary formula identifying the two-cell
boundary with the abelianised commutator product, and (iii) the
comparison-data construction, (iv) chain-map correctness, and (v) the
induced \(H_1\)-isomorphism for the cellular-to-singular comparison.
These are precisely the cellular-chain and Hurewicz-style APIs absent
from Mathlib v4.28.0.}
\lean{JacobianChallenge.Periods.polygon4g_cellular_singular_iso}
\notready
\end{theorem}

\begin{definition}[Orientable typeclass]
\label{def:orientable}
A topological space \(M\) is \emph{orientable} if it carries a
consistent choice of orientation across all charts. (In the current
formalization this is a thin placeholder typeclass with the
real witness deferred to a later round.)
\lean{JacobianChallenge.Periods.Orientable}
\leanok
\end{definition}

\begin{layman}
A surface is \emph{orientable} when you can paint it with a
consistent local sense of clockwise vs counterclockwise — like the
two-coloring of a chessboard, but for ``which way around am I
spinning.'' The Möbius strip and Klein bottle are not orientable;
the sphere, torus, and \(4g\)-gon-with-standard-pattern are. Why
bother? The polygonal-model theorem applies only to orientable
surfaces; the standard relator word \(a_1b_1a_1^{-1}b_1^{-1}\cdots\)
is the orientable normal form. Non-orientable surfaces have a
different normal form (\(a_1^2 a_2^2 \cdots\)), so we need to flag
the orientability hypothesis explicitly.
\end{layman}

\begin{theorem}[Classical input: de Rham theorem]
\label{input:de-rham-theorem}
\uses{lem:deRham-integration-descends,lem:deRham-iso-via-good-cover,lem:deRham-overview-stepwise}
For a smooth manifold \(M\), integration of forms over chains gives
an isomorphism between de Rham and singular cohomology:
\[
  H^{k}_{\mathrm{dR}}(M,\C) \cong H^{k}_{\mathrm{sing}}(M,\C).
\]
\depends{Classical 1931 result of Georges de Rham. Mathlib v4.28.0
has neither de Rham cohomology of manifolds nor the integration
pairing that realises the comparison. The project sketches it in
\code{Jacobian/StageB/DeRhamComplex.lean}
(\(d^2 = 0\), closed/exact subspaces) and
\code{Jacobian/StageB/DeRhamComparison.lean}
(\code{deRham_theorem}, ${\sim}220$ LOC of typed sorries via good
covers + Mayer--Vietoris).}
\lean{JacobianChallenge.StageB.deRham_theorem}
\leanok
\end{theorem}

\begin{layman}
The de Rham theorem is the bridge between two ways of counting
holes in a smooth manifold. The first way uses
\emph{differential forms}: closed forms (\(d\omega = 0\)) modulo
exact ones (\(\omega = d\alpha\)) gives the de Rham cohomology
\(H^k_{\mathrm{dR}}\). The second way uses \emph{singular chains}:
formal sums of continuous simplices, with cycles (\(\partial = 0\))
modulo boundaries, giving singular cohomology. De Rham's theorem
says these are naturally isomorphic, and the isomorphism is
\emph{integration}: given a closed form and a singular cycle, the
integral pairs them. Why bother? It promotes ``how many independent
closed forms are there'' (an analytic question) into ``how many
independent loops are there'' (a topological question), with
matching integer ranks on both sides. Without de Rham, the
analytic side of the genus identity has no topological referent.
\end{layman}

\begin{theorem}[Classical input: Hodge decomposition on a compact Riemann surface]
\label{input:hodge-decomposition}
\uses{input:de-rham-theorem,lem:hodge-laplacian-kernel-finite-dim,lem:hodge-deRham-iso-harmonic,lem:hodge-decomposition-overview-stepwise}
For a compact Kähler manifold \(X\) (in particular a compact Riemann
surface) every de Rham cohomology class has a unique harmonic
representative, and harmonic forms split by bidegree:
\[
  H^{1}_{\mathrm{dR}}(X,\C) \cong
  H^{1,0}_{\bar\partial}(X) \oplus H^{0,1}_{\bar\partial}(X).
\]
\depends{Classical Hodge theorem (W.~V.~D.~Hodge, 1941). Requires
elliptic regularity for the Laplace--Beltrami operator. Mathlib
v4.28.0 has neither the Laplacian on a Riemannian manifold,
elliptic-PDE compactness, harmonic-form decomposition, nor Kähler
\((p,q)\)-bidegree. The project sketches the chain in
\code{Jacobian/StageB/LaplaceBeltrami.lean}
(volume form, Hodge star, codifferential, \(\Delta\) self-adjoint) +
\code{Jacobian/StageB/HarmonicForms.lean}
(\code{hodge_decomposition},
\code{deRhamH_iso_Harmonic}, Poincaré duality) +
\code{Jacobian/StageB/KahlerStructure.lean}
(Dolbeault \((p,q)\)-decomposition, Hodge symmetry).}
\lean{JacobianChallenge.StageB.hodge_decomposition}
\leanok
\end{theorem}

\begin{layman}
The Hodge decomposition theorem says that on a compact Kähler
manifold, every cohomology class has a unique \emph{harmonic}
representative — a form lying in the kernel of the
Laplace--Beltrami operator \(\Delta = d\delta + \delta d\). On a
compact Riemann surface, harmonic 1-forms split further by
\((p,q)\)-bidegree: the holomorphic 1-forms (type \((1,0)\)) and
the anti-holomorphic 1-forms (type \((0,1)\)) each form a
\(g\)-dimensional space, summing to \(2g\). Why bother? This is
the single fact that makes the analytic genus and the topological
genus agree: \(2g_{\mathrm{an}} = \dim H^1_{\mathrm{dR}} =
\mathrm{rank}\, H_1(X,\Z)\) (using de Rham). It is the third
big classical-analysis input on the project's path, alongside
Riemann--Roch and Stokes, and Mathlib has none of the underlying
elliptic-operator + harmonic-form theory yet.
\end{layman}

\begin{theorem}[Classical input: Hodge / de Rham rank for a Riemann surface]
\label{input:hodge-deRham}
\uses{def:topological-genus,input:de-rham-theorem,input:hodge-decomposition,input:dolbeault,lem:hodge-deRham-dim-h1-dR,lem:hodge-deRham-rank-eq-dim,lem:hodge-deRham-overview-stepwise,lem:hodge-deRham-deRham-bridge,lem:hodge-deRham-uct-bridge,lem:hodge-deRham-real-of-complex}
For a compact connected oriented Riemann surface \(X\) of analytic
genus \(g\), the rank of the integer first homology group is
\(2g\):
\[
  \mathrm{rank}_{\Z}\,H_1(X,\Z) = 2 g_{\mathrm{an}}(X).
\]
\depends{The standard route is the Hodge / de Rham theorem:
holomorphic 1-forms together with their conjugates span the
de Rham cohomology \(H^1_{\mathrm{dR}}(X)\), which by de Rham's
theorem is dual to singular \(H_1(X,\R)\); the integer lattice
\(H_1(X,\Z)\) is then full of rank \(2g\). Mathlib v4.28.0 has
neither de Rham cohomology of manifolds, the de Rham theorem, nor
Hodge decomposition for Riemann surfaces. The project ships a
bottom-up sketch of the chain under \code{Jacobian/StageB/}
(\code{DifferentialForms} \(\to\) \code{DeRhamComplex} \(\to\)
\code{DeRhamComparison} \(\to\) \code{LaplaceBeltrami} \(\to\)
\code{HarmonicForms} \(\to\) \code{KahlerStructure} \(\to\)
\code{CoherentSheaves} \(\to\) \code{SerreDuality}; eight files,
${\sim}960$ LOC, all sorry) plus a top-down refinement under
\code{Jacobian/Periods/HodgeDeRham.lean} (30 named sub-leaves).
The project's Lean target is currently a sorry-bound assembly
through which the
\(g_{\mathrm{an}}=g_{\mathrm{top}}\) bridge is wired.}
\lean{JacobianChallenge.Periods.hodge_deRham_rank_eq}
\notready
\end{theorem}

\begin{layman}
The Hodge / de Rham bridge is the deep input that makes the
analytic and topological notions of genus agree at the level of
ranks. Phrased one way: a compact Riemann surface of analytic
genus \(g\) (the dimension of its space of holomorphic 1-forms) has
exactly \(2g\) independent loops in integer homology. Phrased
another way: every closed 1-form is the sum of a holomorphic and
an anti-holomorphic part, and the holomorphic and anti-holomorphic
spaces both have dimension \(g\) — so the de Rham cohomology has
real dimension \(2g\), which by de Rham's theorem matches the
real dimension of \(H_1(X,\R)\), forcing the integer lattice
\(H_1(X,\Z)\) to be free of rank \(2g\). Why bother? Two of the
four anti-hack theorems hinge on this rank computation: the
genus-zero classification (forcing \(g=0\) to mean ``no loops at
all'') and the period-lattice fullness (forcing the
\(2g\) symplectic-basis vectors to span the dual real space).
Mathlib lacks the entire Hodge / de Rham apparatus on manifolds,
making this the third major classical-analysis gap on the path,
alongside Riemann-Roch and Stokes.
\end{layman}

\begin{theorem}[Analytic genus equals topological genus]
\label{thm:analytic-eq-topological-genus}
\uses{thm:polygonal-model,def:topological-genus,input:hodge-deRham}
For a compact connected oriented Riemann surface \(X\),
\[
  g_{\mathrm{an}}(X) = g_{\mathrm{top}}(X).
\]
\depends{Combines the polygonal-model homeomorphism with the
Hodge / de Rham rank computation
(\ref{input:hodge-deRham}): the polygonal model
identifies \(X\) topologically with \(\mathrm{Polygon4g}(g)\) of
known \(H_1\) rank \(2g\), and Hodge / de Rham fixes that rank
to be \(2 g_{\mathrm{an}}(X)\), so dividing by two gives
\(g_{\mathrm{top}}(X) = g_{\mathrm{an}}(X)\). Project assembly:
\code{Jacobian/Periods/AnalyticGenusEqTopologicalGenus.lean}.}
\lean{JacobianChallenge.Periods.analyticGenus_eq_topologicalGenus}
\leanok
\end{theorem}
\begin{proof}\leanok
The Lean proof is a one-line delegate to
\code{stageB_analytic_eq_topological}, itself a sorry-free assembly
of the project's existing
\code{HolomorphicForms.two_analyticGenus_eq_finrank_intH1} (de Rham +
Hodge + UCT chain) plus \code{Nat.mul_div_cancel_left} to identify
\(g_{\mathrm{top}} = \mathrm{rank}_{\Z} H_1 / 2 = (2 g_{\mathrm{an}}) / 2 = g_{\mathrm{an}}\).
\end{proof}

\begin{layman}
The two genera agree. Compute the topological genus of \(X\) by
homology; compute the analytic genus as the dimension of
holomorphic 1-forms; they're the same number. The proof factors
through the polygonal model: the homeomorphism \(X\cong
\mathrm{Polygon4g}(g)\) lets us compute \(H_1(X)\) on the
combinatorially explicit polygon (where the symplectic loops are
visible), giving rank \(2g\), so the topological genus is \(g\).
That \(g\) was put in by hand to match the analytic genus, closing
the loop. Why bother? It's the bridge that makes the
genus-zero classification (\ref{thm:genus-zero-classification})
work: ``analytic genus \(=0\) means topological genus \(=0\) means
\(H_1=0\) means homeomorphic to \(S^2\).''
\end{layman}

\begin{proofsketch}
Apply Hodge / de Rham to compute \(\mathrm{rank}_\Z H_1(X,\Z) = 2 g_{\mathrm{an}}(X)\),
then divide by two. The Lean target
\code{JacobianChallenge.Periods.analyticGenus\_eq\_topologicalGenus}
is a sorry-free assembly in
\code{Jacobian/Periods/PeriodFunctional.lean}, delegating to two
named sub-obligations: \code{hodge\_deRham\_rank\_eq} (the analytic
core: \(2g_{\mathrm{an}} = \mathrm{rank}_\Z H_1\)), itself
decomposed into a multi-file frontier-obligation tree under
\code{Jacobian/HolomorphicForms/HodgeDeRhamRank.lean} and
\code{Jacobian/Periods/HodgeDeRham.lean}; and the canonical
\(\mathrm{rank}_\Z H_1 / 2\) definition of \code{topologicalGenus}
from \code{Jacobian/Periods/TopologicalGenus.lean}. The two-line
assembly body is \code{omega} after rewriting through
\code{hodge\_deRham\_rank\_eq}.
\end{proofsketch}

\subsection{Analysis on the polygon}

The interior of the standard fundamental polygon is, by construction,
homeomorphic to an open disk: the side-pairing only identifies
\emph{boundary} points. This subsection records the consequences
needed by the next chapter: that the interior is open in the
quotient, that the quotient map is injective there, and that every
holomorphic 1-form on it admits a primitive.

\begin{lemma}[Quotient is injective on the polygon interior]
\label{lem:mk-injOn-openDisk}
\uses{def:polygon4g}
The polygon-quotient map \(\mathrm{mk}_g\) is injective on the open
unit disk: distinct interior points stay distinct after taking
equivalence classes.
\lean{JacobianChallenge.Periods.PrimitiveOnPolygon.mk_injOn_openDisk}
\leanok
\end{lemma}
\begin{proof}\leanok
Every generating pair in the side-pairing relation has both points on
the boundary circle. Hence the equivalence closure cannot relate any
strictly interior point to any other point.
\end{proof}

\begin{layman}
Inside the open disk \(\{|z|<1\}\), the quotient map is one-to-one.
The side-pairing relation only ever glues boundary points
(\(|z|=1\)) to other boundary points; it never touches the interior.
Why bother? It says the interior of the polygon embeds verbatim
into the quotient — making it possible to talk about ``the inside
of the polygon'' as a single open subset on which classical
complex-analysis lemmas apply directly.
\end{layman}

\begin{lemma}[Polygon interior is open in the quotient]
\label{lem:mk-image-openDisk-isOpen}
\uses{def:polygon4g,lem:mk-injOn-openDisk}
The image of the open unit disk under \(\mathrm{mk}_g\) is open in
\(\mathrm{Polygon4g}(g)\).
\lean{JacobianChallenge.Periods.PrimitiveOnPolygon.mk_image_openDisk_isOpen}
\leanok
\end{lemma}
\begin{proof}\leanok\end{proof}

\begin{layman}
The image of the open disk in the quotient is itself an open subset
of the polygon. The proof goes through the quotient topology:
the preimage of the image equals the open disk again (because the
side-pairing only identifies boundary points to boundary points), and
the open disk is open in the original disk. Why bother? Stokes,
Cauchy, and primitive arguments need to know the polygon's
interior is open so that ``differentiable on the open polygon''
is well-defined.
\end{layman}

\begin{lemma}[Holomorphic primitive on the open disk]
\label{lem:holomorphic-primitive-openDisk}
Every holomorphic function on the open unit disk admits a holomorphic
primitive on it.
\depends{Cauchy's theorem on a convex domain (Mathlib has this in
\code{Mathlib.Analysis.Complex.HasPrimitives}).}
\lean{JacobianChallenge.Periods.PrimitiveOnPolygon.holomorphic_has_primitive_openDisk}
\leanok
\end{lemma}
\begin{proof}\leanok\end{proof}

\begin{layman}
On a convex region (and the open unit disk is convex), every
holomorphic function has a primitive: a holomorphic antiderivative.
The proof in Mathlib goes via Cauchy's theorem on contractible
domains. Why bother? The next-chapter Stokes-on-the-polygon
argument turns ``integrate \(\omega\) along a polygon edge'' into
``evaluate the primitive \(F\) at the endpoints,'' so we need
\(F\) to exist.
\end{layman}

\begin{lemma}[Primitive of a holomorphic 1-form on the polygon]
\label{lem:primitive-on-polygon}
\uses{thm:polygonal-model,lem:mk-injOn-openDisk,lem:mk-image-openDisk-isOpen,lem:holomorphic-primitive-openDisk}
On the interior of the standard fundamental polygon, every
holomorphic 1-form \(\omega\) admits a holomorphic primitive
\(F\) with \(dF=\omega\).
\lean{JacobianChallenge.Blueprint.primitive_on_polygon}
\leanok
\end{lemma}
\begin{proof}\leanok\end{proof}

\begin{layman}
Pull the holomorphic 1-form back to the polygon's flat interior.
The interior is simply connected (in fact convex, since it's the
open unit disk). On any simply connected flat region, every
closed 1-form admits a primitive — a single function whose
differential equals the form. Compare: on the punctured plane
\(d\theta\) has no global primitive because of the puncture-loop;
on the disk it does. Apply this to \(\omega\); the resulting
primitive \(F\) is the bookkeeping device that turns wedge
integrals over the surface into boundary integrals around the
polygon. Why bother? Without a primitive, the next chapter's
bilinear-from-Stokes argument cannot fold edges together pairwise
to pick up periods. With it, the entire Riemann bilinear
identity falls out by an integration by parts.
\end{layman}

\begin{proofsketch}
The Lean target in
\code{Jacobian/Blueprint/Sec03/PrimitiveOnPolygon.lean} is now
\textbf{fully sorry-free}. The headline
\code{primitive\_on\_polygon} is a substantive Prop on the
project-internal stand-in
\code{chartPullbackPrimitive}\,\(g\;h\): for every chart-pullback
\(h:\C\to\C\) with \(\code{DifferentiableOn}\;\C\;h\;\mathrm{OpenDisk}\),
the chart-roundtrip data exists — a holomorphic primitive on the
disk together with the polygon chart's
\code{InjOn}/\code{IsOpen}-image properties. The four sub-leaves
(\code{primitive\_on\_polygon\_mk\_injOn},
\code{primitive\_on\_polygon\_image\_isOpen},
\code{primitive\_on\_polygon\_disk\_primitive}, and the former frontier
obligation \code{primitive\_on\_polygon\_chart\_pullback} — now
lifted to "continuous functions on \(\mathrm{Polygon4g}(g)\) pull
back to continuous functions on the open disk" and refined across
5 passes through \code{primitive\_on\_polygon\_mk\_continuous}
(\code{continuous\_quotient\_mk'}),
\code{primitive\_on\_polygon\_chart\_continuous}
(\code{Continuous.comp}), and the holomorphic-upgrade shadow
\code{primitive\_on\_polygon\_chart\_pullback\_differentiable}
together with \code{openDisk\_subset\_closedBall}
(\code{Metric.ball\_subset\_closedBall})) are all sorry-free; the
headline is a sorry-free assembly. The eventual lift to
\(\code{HolomorphicPrimitive}\) on \(\mathrm{Polygon4g}(g)\) reduces
to this stand-in once the complex-manifold structure on
\(\mathrm{Polygon4g}\) is in place.
\end{proofsketch}

\begin{formalizationnote}
The lemma now has a sorry-free chart-pullback-level Lean proof
(\ref{lem:mk-injOn-openDisk}, \ref{lem:mk-image-openDisk-isOpen},
\ref{lem:holomorphic-primitive-openDisk} feed the headline through
\code{chartPullbackPrimitive}). The remaining infrastructure step
is to upgrade the chart-pullback shadow to a manifold-side
\(\code{HolomorphicOneForm}\) pullback once
\(\mathrm{Polygon4g}\) carries a complex-manifold structure;
\code{primitive\_on\_polygon\_chart\_pullback\_differentiable} in
\code{Sec03/PrimitiveOnPolygon.lean} tracks that obligation.
\end{formalizationnote}
